% foundations/aim_roadmap.tex

\chapter*{AIM Roadmap}
\addcontentsline{toc}{chapter}{AIM Roadmap}

\label{chap:aim-roadmap}

The present thesis develops Alignment-Based Information Modelling (AIM)
through a sequence of increasingly richer modelling concepts. The roadmap
below provides a global overview before the individual concepts are
introduced in detail.

\begin{figure}[ht]
\centering

\begin{tikzpicture}[
    node distance=9mm,
    every node/.style={
        draw,
        rounded corners,
        align=center,
        minimum width=45mm,
        minimum height=8mm
    },
    >=Stealth
]

\node (classic)
{
Classical Alignment\\
Engineering
};

\node (curvature) [below=of classic]
{
Curvature Space
};

\node (pose2) [below=of curvature]
{
pose2
};

\node (pose3) [below=of pose2]
{
pose3
};

\node (vehicle) [below=of pose3]
{
Vehicle Space
};

\node (runtime) [below=of vehicle]
{
Runtime State
};

\node (realization) [below=of runtime]
{
Realization
};

\node (optimization) [below=of realization]
{
Optimization
};

\node (aim) [below=of optimization]
{
Alignment-Based\\
Information Modelling (AIM)
};

\node (pim) [below=of aim]
{
Process Information\\
Modelling (PIM)
};

\draw[->] (classic) -- (curvature);
\draw[->] (curvature) -- (pose2);
\draw[->] (pose2) -- (pose3);
\draw[->] (pose3) -- (vehicle);
\draw[->] (vehicle) -- (runtime);
\draw[->] (runtime) -- (realization);
\draw[->] (realization) -- (optimization);
\draw[->] (optimization) -- (aim);
\draw[->] (aim) -- (pim);

\end{tikzpicture}

\caption{
Conceptual roadmap of the AIM framework.
The thesis proceeds from intrinsic railway geometry toward runtime,
realization, and optimization concepts, culminating in
Alignment-Based Information Modelling (AIM) as part of a broader
Process Information Modelling (PIM) perspective.
}
\label{fig:aim-roadmap}

\end{figure}

\section*{Reading Guide}

The development of the thesis follows four major stages.

\begin{enumerate}

\item
\textbf{Geometry.}
Railway alignments are reinterpreted as curvature-driven objects. The
central geometric state is introduced as \emph{pose2}, which emerges
from curvature integration rather than from explicit coordinate
construction.

\item
\textbf{State.}
The geometric description is extended to spatial and operational
contexts through \emph{pose3}, vehicle-space concepts, and runtime
state representations.

\item
\textbf{Reality and Runtime.}
Real railway systems are never ideal. Measurement uncertainty,
realization effects, construction tolerances, and operational
influences require runtime-aware modelling concepts.

\item
\textbf{Optimization.}
Once geometry, state, runtime, and realization are unified, railway
engineering can be reformulated as an optimization problem in state
space rather than purely as a geometric construction problem.

\end{enumerate}

The thesis begins with railway geometry, extends geometry into state,
extends state into runtime and realization, and finally arrives at
optimization-oriented Alignment-Based Information Modelling (AIM). Within
the broader perspective proposed here, AIM may be interpreted as one
specialization of a more general Process Information Modelling (PIM)
framework.
